\section{Einleitung}
\label{sec:Einleitung}

Das Verhalten von realen Gasen spielt in der experimentellen Thermodynamik eine wichtige Rolle.
Im Gegensatz zu idealen Gasen, die durch die ideale Gasgleichung beschrieben werden können,
zeigen reale Gase Abweichungen von diesem Modell, insbesondere bei hohen Drücken und niedrigen Temperaturen.
Besonders deutlich treten die Abweichungen vom idealen Verhalten
im Bereich des Phasenübergangs zwischen Gas- und Flüssigphase auf.
Unterhalb einer charakteristischen Temperatur koexistieren beide
Phasen bei gleichem Druck, wobei das Volumen beim Übergang sprunghaft
wechselt. Mit steigender Temperatur nähert sich diese Phasengrenze
einem kritischen Punkt an, an dem sich Gas- und Flüssigphase nicht
mehr unterscheiden lassen. Am kritischen Punkt verschwinden die
Dichteunterschiede zwischen beiden Phasen sowie die Grenzflächenspannung.
Das Studium dieses Punktes ist von besonderer physikalischer
Bedeutung, da hier kollektive molekulare Effekte dominieren und
klassische Näherungen an ihre Grenzen stoßen.
Um das Verhalten realer Gase besser zu verstehen, wurden verschiedene Modelle entwickelt, darunter die Van-der-Waals-Gleichung,
die Korrekturen für diese Effekte einführt.
Die Van-der-Waals-Gleichung beschreibt zwar qualitativ das Auftreten
eines kritischen Punktes, liefert jedoch unterhalb der kritischen
Temperatur Isothermen mit einem oszillierenden Verlauf. In einem
bestimmten Volumenbereich ergibt sich dabei eine negative Steigung,
was einer negativen Kompressibilität entspricht und thermodynamisch
instabil ist. Dieser Bereich ist daher unphysikalisch und wird
experimentell nicht beobachtet.
Zur Korrektur dieses Problems wird die sogenannte Maxwell-Konstruktion
verwendet. Dabei wird der instabile Abschnitt durch eine horizontale
Gerade ersetzt, die so gewählt wird, dass die Flächen ober- und
unterhalb der Geraden gleich groß sind. Diese Maxwell-Gerade beschreibt
den realen Phasenübergang bei konstantem Druck und stellt sicher,
dass die freien Energien beider Phasen identisch sind.
In diesem Experiment wird das Verhalten von Schwefelhexafluorid (\(SF_6\)) untersucht, indem der Druck in Abhängigkeit vom Volumen bei verschiedenen
Temperaturen gemessen wird und die Ergebnisse mit den Vorhersagen der Van-der-Waals-Gleichung verglichen werden.
Schwefelhexafluorid (SF$_6$) eignet sich besonders gut für die
experimentelle Untersuchung realer Gase, da seine kritischen
Parameter bei vergleichsweise moderaten Temperaturen und Drücken
liegen. Dadurch kann der kritische Punkt im Labormaßstab sicher
und mit vertretbarem technischem Aufwand erreicht werden.
Zudem ist SF$_6$ chemisch stabil, ungiftig und nicht brennbar,
was eine kontrollierte Durchführung des Experiments ermöglicht.